\chapter{Complexity — Karp Lipton Search Family}
%%% generated by the blueprint pipeline from TCSlib.Complexity.CircuitComplexity.KarpLiptonSearchFamily
\section{Overview}

Declarations of the Lean module \texttt{TCSlib.Complexity.CircuitComplexity.KarpLiptonSearchFamily}.

\section{Declarations}

\begin{theorem}[Acyclicity of two constant gates]
\lean{BoolCircuit.gatesAcyclic_consts}
\label{BoolCircuit.gatesAcyclic_consts}
\leanok
\uses{BoolCircuit.GatesAcyclic, BoolCircuit.constGate}
\difficulty{4}
For every natural number $k$, the list consisting of a false constant gate followed by a
true constant gate is acyclic starting from index bound $k$.
\end{theorem}

\begin{definition}[Circuit rewiring with two constant inputs]
\lean{BoolCircuit.DAGCircuit.rewire}
\label{BoolCircuit.DAGCircuit.rewire}
\leanok
\uses{BoolCircuit.DAGCircuit, BoolCircuit.GatesAcyclic.append, BoolCircuit.constGate, BoolCircuit.embedGates, BoolCircuit.gatesAcyclic_consts, BoolCircuit.gatesAcyclic_embedGates}
Given a circuit with $N$ input vertices, a natural number $k$, and a list of $N$ source
vertices, suppose the list has length $N$ and every listed vertex is less than $k+2$.
Rewiring produces a circuit on $k$ inputs by placing constant-false and constant-true
gates first, then appending the original gates with their inputs redirected according to
the source list; its output is the vertex $k+2$ plus the number of original gates.
\end{definition}

\begin{theorem}[Evaluation After Rewiring]
\lean{BoolCircuit.DAGCircuit.rewire_eval}
\label{BoolCircuit.DAGCircuit.rewire_eval}
\leanok
\uses{BoolCircuit.DAGCircuit.eval, BoolCircuit.DAGCircuit.rewire, BoolCircuit.DAGCircuit.values, BoolCircuit.DAGGate.eval, BoolCircuit.constGate, BoolCircuit.runWith, BoolCircuit.runWith_append, BoolCircuit.runWith_embedGates_getD}
\difficulty{4}
Let $N$ and $k$ be nonnegative integers, let $C$ be a Boolean circuit with $N$ inputs,
and let $\text{srcs}$ be a list of $N$ vertex indices, each less than $k+2$. The circuit
$C$ is assumed to have gates in topological order, each reading only earlier vertices,
and a valid output vertex. For every assignment $x$ to $k$ bits, the rewired circuit's
output equals the output of $C$ when its $i$th input is taken from the list consisting
of $x$ followed by $0,1$, at index $\text{srcs}[i]$.
\end{theorem}

\begin{theorem}[Rewiring Preserves Fan-In Two]
\lean{BoolCircuit.DAGCircuit.rewire_isFaninTwo}
\label{BoolCircuit.DAGCircuit.rewire_isFaninTwo}
\leanok
\uses{BoolCircuit.DAGCircuit.IsFaninTwo, BoolCircuit.DAGCircuit.rewire, BoolCircuit.constGate, BoolCircuit.faninTwo_embedGates}
\difficulty{4}
Let $C$ be a Boolean circuit with $N$ inputs, and let $k$ be a nonnegative integer.
Suppose the source list has length $N$, each source index is less than $k+2$, and every
gate of $C$ has distinct input vertices, every NOT gate has exactly one input, and every
gate has at most two inputs. Then the circuit obtained by rewiring $C$ using those
sources, with two additional constant inputs $0$ and $1$, has distinct input vertices at
every gate, exactly one input at each NOT gate, and at most two inputs at every gate.
\end{theorem}

\begin{theorem}[Size after rewiring]
\lean{BoolCircuit.DAGCircuit.rewire_size}
\label{BoolCircuit.DAGCircuit.rewire_size}
\leanok
\uses{BoolCircuit.DAGCircuit.rewire, BoolCircuit.DAGCircuit.size}
\difficulty{3}
Let $C$ be a finite acyclic Boolean circuit with $N$ input vertices and gate list $G$,
and let $k$ be a nonnegative integer. Suppose a list of $N$ source vertices has every
entry less than $k+2$. Rewiring $C$ using these sources to obtain a circuit with $k$
inputs gives a circuit of size $k+3+|G|$.
\end{theorem}

\begin{definition}[Query source list]
\lean{BoolCircuit.KarpLipton.querySrcs}
\label{BoolCircuit.KarpLipton.querySrcs}
\leanok
For natural numbers $n,j,m$, form a list by concatenating two copies of each number from
$0$ through $n-1$, the two entries $n+j$ and $n+j+1$, then $m-j$ copies of $n+j+1$ (zero
copies if $m<j$), followed by $n+j$ and the list $n,n+1,\ldots,n+j-1$.
\end{definition}

\begin{theorem}[Length of the query source list]
\lean{BoolCircuit.KarpLipton.length_querySrcs}
\label{BoolCircuit.KarpLipton.length_querySrcs}
\leanok
\uses{BoolCircuit.KarpLipton.querySrcs}
\difficulty{3}
For natural numbers $n,j,m$, the list of query sources has length $2n+2+(m-j)+1+j$,
where $m-j$ denotes truncated subtraction.
\end{theorem}

\begin{theorem}[Bound on query source indices]
\lean{BoolCircuit.KarpLipton.querySrcs_lt}
\label{BoolCircuit.KarpLipton.querySrcs_lt}
\leanok
\uses{BoolCircuit.KarpLipton.querySrcs}
\difficulty{3}
For natural numbers $n,j,m$, every entry of the query source list is less than $n+j+2$.
\end{theorem}

\begin{theorem}[Query source encoding]
\lean{BoolCircuit.KarpLipton.querySrcs_map}
\label{BoolCircuit.KarpLipton.querySrcs_map}
\leanok
\uses{BoolCircuit.KarpLipton.querySrcs, Turing.pairEncode}
\difficulty{5}
Let $n,j,m$ be natural numbers, and let $y$ and $p$ be bit strings with lengths $n$ and
$j$, respectively. Replace each query source index by the bit at that position in the
concatenation of $y$, $p$, and $01$, using $0$ when the index is out of range. The
resulting bit string is the self-delimiting pairing of $y$ with the string consisting of
$m-j$ copies of $1$, followed by $0$, followed by $p$.
\end{theorem}

\begin{theorem}[Circuit-family evaluation and language membership]
\lean{BoolCircuit.family_eval_getD}
\label{BoolCircuit.family_eval_getD}
\leanok
\uses{BoolCircuit.DAGCircuit.eval, BoolCircuit.DAGCircuitFamily, BoolCircuit.DAGCircuitFamily.language, BoolCircuit.DAGCircuitFamily.mem_language_iff}
\difficulty{3}
Let $F$ be a family of directed acyclic Boolean circuits, let $N$ be a natural number,
and let $q$ be a Boolean string of length $N$. The circuit in the family for inputs of
length $N$, evaluated on the length-$N$ input whose $i$th bit is the $i$th bit of $q$
(or false when that index is out of range), outputs true if and only if $q$ belongs to
the language recognized by the family.
\end{theorem}

\begin{theorem}[Search circuits from polynomial circuit complexity]
\lean{BoolCircuit.exists_searchCircuit_of_NP_subset_PPoly}
\label{BoolCircuit.exists_searchCircuit_of_NP_subset_PPoly}
\leanok
\uses{BoolCircuit.DAGCircuit, BoolCircuit.DAGCircuit.IsFaninTwo, BoolCircuit.DAGCircuit.eval, BoolCircuit.DAGCircuit.rewire, BoolCircuit.DAGCircuit.rewire_eval, BoolCircuit.DAGCircuit.rewire_isFaninTwo, BoolCircuit.DAGCircuit.rewire_size, BoolCircuit.DAGCircuit.size, BoolCircuit.Extends, BoolCircuit.KarpLipton.length_querySrcs, BoolCircuit.KarpLipton.querySrcs, BoolCircuit.KarpLipton.querySrcs_lt, BoolCircuit.KarpLipton.querySrcs_map, BoolCircuit.MultiDAGCircuit, BoolCircuit.MultiDAGCircuit.IsFaninTwo, BoolCircuit.MultiDAGCircuit.eval, BoolCircuit.MultiDAGCircuit.size, BoolCircuit.PPoly, BoolCircuit.constCircuit, BoolCircuit.constCircuit_isFaninTwo, BoolCircuit.family_eval_getD, BoolCircuit.getD_map_getD, BoolCircuit.searchCircuit, BoolCircuit.searchCircuit_correct, BoolCircuit.searchCircuit_isFaninTwo, BoolCircuit.searchCircuit_size_le, BoolCircuit.searchState_struct, Complexity.KarpLipton.mem_prefixLangOf_iff, Complexity.KarpLipton.prefixLangOf_mem_NP, Complexity.NP, Complexity.P, Complexity.PolyHierarchy.UnaryPT, Complexity.PolyHierarchy.unaryPT_poly, Complexity.polyTimeComputable_id, Turing.length_pairEncode, Turing.pairEncode}
\difficulty{7}
Assume every language in $mathsf{NP}$ has polynomial-size Boolean circuits, and let $V$
be a language in $mathsf{P}$. For any natural numbers $C$ and $c$, there are natural
numbers $q_0$ and $e$ such that, for every input length $n$, there is a fan-in-two
circuit with $C(n+1)^c$ outputs and size at most $q_0(n+1)^e$ that preserves witnesses
for $V$: whenever some string $v$ of length $C(n+1)^c$ makes the encoded pair of an
input $x$ and $v$ belong to $V$, the circuit's output on $x$ also makes that encoded
pair belong to $V$.
\end{theorem}
